\section{Quasiprimitive extraction and the minimal-source alternating cutoff}
\label{sec:quasiprimitive-extraction}

\paragraph{Accepted status and provenance.}
D-TD62, question Q2fr, is written by
\texttt{/root/sixfolds}. Root proposed checking whether homogeneity,
rather than primitivity, suffices for the complete extraction and
cutoff argument. The residual-multiplicity proof below is new;
\texttt{/root/imprimitive} independently reconstructed that specific
proof as a targeted check before this full section was written.
The original-complete-layer replacement for the risky complement
step was developed in discussion with root. The accepted character,
actual-cover and natural-polarization inputs retain their prior
source status. The author has read all 688 initial proof lines, including
every displayed argument and each downstream interface, and records PASS.
Complete root and independent whole-proof reviews also pass. No numerical
execution is needed or claimed.

\subsection*{1. The exact homogeneity hypothesis}

An irreducible complex representation of a finite group is
\emph{quasiprimitive} if its restriction to every normal subgroup
is homogeneous: it is a direct sum of copies of one irreducible
representation of that subgroup. A finite projective image is
called quasiprimitive when its finite special-linear inverse image
has this property in the displayed representation.

This property is unchanged by adjoining finite scalar matrices,
or by passing between finite linear lifts of the same projective
image. Here is the normal-subgroup verification. If
\(E=JZ\), where \(Z\) is a finite central scalar group and \(J\)
acts quasiprimitively, then for \(D\triangleleft E\) one has
\(DZ=(DZ\cap J)Z\), with \(DZ\cap J\triangleleft J\).
Adjoining scalar matrices does not change an invariant subspace
or an isotypic decomposition. Thus restriction to \(DZ\), and
then to \(D\), is homogeneous. Conversely a normal subgroup of
\(J\) becomes normal in \(JZ\) after adjoining \(Z\), so the
same argument restricts the property back to \(J\).
Two finite lifts have a common finite scalar enlargement:
the scalars needed to adjust their finitely many representatives
have bounded root-of-unity orders. This proves the assertion
about lifts. Pullback along a surjection also preserves the
property, since the image of a normal subgroup is normal and
the kernel acts trivially.

No assertion is made that quasiprimitivity is the same as
primitivity, or that an arbitrary tensor factor of a quasiprimitive
representation must be quasiprimitive.

\subsection*{2. The residual multiplicity quotient}

Let \(P\subset\operatorname{PGL}(V)\) be finite and
quasiprimitive, with \(d=\dim V>1\). Choose
\(1\ne N\triangleleft P\). Let \(J\) be the finite special-linear
inverse image of \(P\), let \(Z\) be its scalar kernel, and let
\(R\triangleleft J\) be the inverse image of \(N\).
Homogeneity and the usual isotypic evaluation give
\[
 V=W\otimes U,\qquad
 R|_V=\rho_W(R)\otimes I_U,\qquad
 U=\operatorname{Hom}_R(W,V).
 \tag{QPX.1}
\]
The representation of \(R\) on \(W\) is faithful: its kernel
would act trivially on \(V\), whereas \(J\) is an actual matrix
group. If \(\dim W=1\), then \(R\) is scalar on \(V\), contrary
to \(N\ne1\). Hence \(\dim W\ge2\).

The matrix algebra generated by \(R\) is
\(\operatorname{End}(W)\otimes I\), and its centralizer is
\(I\otimes\operatorname{End}(U)\). The complete matrix-unit
argument of D-TD11 therefore applies verbatim after replacing
its homogeneous-restriction premise by the present hypothesis.
It gives genuine finite projective homomorphisms
\[
 \pi_W:P\longrightarrow\operatorname{PGL}(W),\qquad
 \pi_U:P\longrightarrow\operatorname{PGL}(U),\qquad
 g=A_g\otimes B_g.
 \tag{QPX.2}
\]
Only the projective classes of the two factors are chosen;
their ambiguity consists of inverse scalar changes. Their
multiplicativity follows from uniqueness of these projective
classes. Both factor actions are irreducible, since an invariant
subspace tensored with the other whole factor would contradict
irreducibility of \(V\).

\begin{lemma}[Residual homogeneity]
If \(\dim U>1\), the actual finite projective image \(P_U\)
on \(U\) is quasiprimitive. It has dimension at most \(d/2\)
and is nontrivial.
\end{lemma}

\begin{proof}
Let \(J_U\subset\operatorname{SL}(U)\) be the finite inverse
image of \(P_U\), and form the finite group
\[
 \widetilde J
 =\{(g,b)\in J\times J_U:\pi_U([g])=[b]\},
 \qquad
 \widehat R=\{(r,I):r\in R\}.
 \tag{QPX.3}
\]
Both projections are onto. The representation of
\(\widetilde J\) on \(V\) through the first projection is
quasiprimitive by Section~1. For every pair \((g,b)\) there
is a unique matrix \(A(g,b)\) such that
\(g=A(g,b)\otimes b\). Uniqueness and multiplication in
\(\widetilde J\) show that \(A\) is a genuine linear
representation. On the normal subgroup \(\widehat R\), it
is precisely the irreducible representation \(\rho_W\).

Take any \(S\triangleleft J_U\), and let
\(M=\operatorname{pr}_2^{-1}(S)\triangleleft\widetilde J\).
Then \(M\) contains \(\widehat R\) and maps onto \(S\).
Write \(U|_S=\bigoplus_i a_i U_i\), with \(U_i\) pairwise
inequivalent irreducibles. Each \(W\otimes U_i\), with the
action \(A(g,b)\otimes U_i(b)\), is an \(M\)-module.
Schur's lemma on \(\widehat R\) gives the exact equality
\[
 \operatorname{Hom}_M(W\otimes U_i,W\otimes U_j)
       =I_W\otimes\operatorname{Hom}_S(U_i,U_j).
 \tag{QPX.4}
\]
Indeed a map commuting with \(\widehat R\) is \(I_W\otimes f\);
commuting with \(A(g,b)\otimes U_i(b)\) is then exactly the
\(S\)-intertwining condition, since \(M\to S\) is onto.
Complete reducibility for the finite group \(M\) shows that
these are irreducible, pairwise inequivalent \(M\)-modules.
Quasiprimitivity of the pulled-back \(V\) forces only one
index \(i\). Thus \(U|_S\) is homogeneous for every normal
\(S\), as required. The dimension bound follows from
\(\dim W\ge2\). A trivial projective action on a space of
dimension greater than one is reducible, proving nontriviality.
\end{proof}

\subsection*{3. Minimal cores and the Heisenberg block argument}

Now choose \(N\) minimal among the nontrivial normal subgroups
of the current quasiprimitive projective image. The elementary
minimal-normal argument in D-TD12 gives either an elementary
abelian \(p\)-group or a direct product \(S^t\) of isomorphic
nonabelian simple groups; in the latter case the current group
permutes the \(t\) factors transitively.

We spell out why its core identification does not use primitivity.
In the abelian case, commutators of lifts define an alternating
pairing \(N\times N\to\mu_p\). Its radical is normal in the
current projective group, and hence is zero or all of \(N\).
The latter case makes \(R\) abelian and its irreducible module
\(W\) one-dimensional, a contradiction. Thus the pairing is
nondegenerate. The trace orthogonality argument of D-TD12
then gives
\[
 Q=N\simeq\mathbf F_p^{2h},\qquad
 a=\dim W=p^h,\qquad |Q|=a^2.
 \tag{QPX.5}
\]
This is an actual irreducible projective Heisenberg module,
without any splitting of its normalizer.

In the nonabelian case take the inverse images of the simple
factors in \(R\), and their derived groups \(E_i\).
The complete central-extension argument in D-TD12 makes each
\(E_i\) perfect and quasisimple. Distinct \(E_i\) commute:
their commutators are central, and a commutator map from a
perfect group to an abelian group is trivial. The same proof
gives \(R=Z E_1\cdots E_t\). Pulling the irreducible
\(R\)-module back to this central product gives
\[
 W=\bigotimes_{i=1}^t W_i,\qquad
 E_i\text{ faithful and irreducible on }W_i,\qquad
 \dim W_i\ge2.
 \tag{QPX.6}
\]
Every fixed-component local operator normalizes its represented
perfect group honestly: projective normalization becomes ordinary
normalization on taking derived groups. These are exactly the
core and full-matrix-algebra conclusions of D-TD12.

For clarity, the coarse image \(P_W=\pi_W(P)\) has \(N\)
as an injectively represented minimal normal subgroup.
If an element of \(R\) were scalar on \(W\), it would be
scalar on \(V\), so its class in \(N\) is trivial.
Minimal normality therefore passes through \(\pi_W\).
In the Heisenberg case the projective centralizer of \(Q\)
is \(Q\), by the complete Weyl-basis calculation D-TD05.
In the nonabelian case the projective centralizer of \(N\)
is trivial. Namely, the scalar commutator with the perfect
central product \(E_1\cdots E_t\) is a linear character,
so is trivial, after which irreducibility and Schur's lemma
give only scalar matrices.

If \(D\triangleleft P_W\) has \(D\cap N=1\), then
\([D,N]=1\). These two centralizer descriptions force \(D=1\).
Thus every nontrivial normal subgroup of \(P_W\) contains
the irreducibly represented \(N\). Adjoining the scalar
kernel shows that its lift acts irreducibly; the same is
true before adjoining scalars. Normal subgroups with trivial
projective image are scalar. Consequently this coarse image
is itself quasiprimitive.

\begin{lemma}[Primitivity of an extracted Heisenberg image]
The coarse Heisenberg image \(P_W\) just constructed is primitive.
Its actual action on the original \(\mathbf F_p\)-space \(Q\)
is irreducible.
\end{lemma}

\begin{proof}
Any invariant \(\mathbf F_p\)-subspace of \(Q\) is a normal
subgroup of \(P_W\), so minimal normality proves the last
assertion. Suppose that \(W\) has a nontrivial direct-sum
system of blocks permuted by \(P_W\). The orbit sum of
any \(Q\)-orbit of blocks is invariant under the irreducible
linear inverse image of \(Q\). Therefore \(Q\) is transitive
on these blocks. Its permutation kernel is a \(P_W\)-invariant
subspace of \(Q\), and is either \(Q\) or zero. The first
case leaves each block \(Q\)-invariant, contradicting
irreducibility. In the second case a faithful transitive
action of the abelian group \(Q\) is regular: all point
stabilizers coincide and equal the permutation kernel.
There would be \(|Q|=a^2\) nonzero blocks in a vector space
of dimension \(a>1\), which is impossible.
\end{proof}

\subsection*{4. Complete recursion and complements from that recursion}

The residual lemma permits iteration on \(P_U\).
Its dimension drops by at least a factor two at each step.
Consequently the D-TD12 construction, with the homogeneity
and residual steps above, gives a complete factorization
\[
 V=\bigotimes_{j=1}^{s} V_j,\qquad
 \dim V_j\ge2,\qquad s\le\lfloor\log_2\dim V\rfloor.
 \tag{QPX.7}
\]
The refined pieces are faithful irreducible quasisimple modules
or irreducible projective Heisenberg modules. Their embedded
matrix algebras are permuted by the original group. At a fixed
quasisimple algebra the local matrix normalizes its represented
perfect group. No residual multiplicity space is discarded.
All later images are actual projective quotients pulled back
through the preceding quotient maps.

Keep also the \emph{coarse layers}: one whole \(W\) from
each extraction step, before splitting a nonabelian layer into
the \(W_i\). Each coarse layer is separately stable, and its
image \(P_j\) contains the just-described minimal normal core.
The individual refined factors in a nonabelian layer can be
permuted. No assertion is made that a later core lifts as a
subgroup of the original coefficient group.

For any subcollection \(\mathscr D\) of whole coarse layers,
let \(U_{\mathscr D}\) be their tensor product and \(P_{\mathscr D}\)
its actual projective image. There is a canonical subdirect
embedding
\[
       P_{\mathscr D}\lhook\joinrel\longrightarrow
                       \prod_{j\in\mathscr D}P_j.
 \tag{QPX.8}
\]
Indeed an original element scalar on \(U_{\mathscr D}\)
conjugates each of its stable full coarse matrix algebras
trivially, and thus is scalar on each coarse factor.
Conversely scalar local matrices have scalar tensor product.
The same reasoning applies to the ratio of two original
elements, proving well-definedness and injectivity of the
map. Each coordinate is onto by definition of the actual
images. This is the exact subcollection argument in D-TD17.
It is not a new recursive decomposition of \(U_{\mathscr D}\).

Let \(m_{\mathscr D}\) be the largest alternating natural
degree among the quasisimple refined pieces in these coarse
layers, or \(1\) if none occurs. The full HD-FT20--21 and
D-TD08 normalizer bounds apply to (QPX.7):
\[
 \frac{\nu(S)}{\dim U_{\mathscr D}}
 \ge \frac1{\max\{80,m_{\mathscr D}-1\}}
 \quad\text{for every nonscalar }S\in P_{\mathscr D}.
 \tag{QPX.9}
\]
Here \(\nu\) means the codimension of the largest eigenspace,
so the bound is independent of scalar representatives.
The local floors are \(1/80\) for nonalternating quasisimple
cores, \(1/(m-1)\) for alternating degree \(m\ge7\), \(1/46\)
for the smaller alternating covers, and at least \(1/4\)
for Heisenberg normalizers. HD-FT20 expressly allows faithful
irreducible quasisimple modules without a primitivity premise.
An operator moving a refined factor has the HD-FT18 quarter
floor. An operator fixing all such factors has the minimum
local floor by spectral convolution. This verifies every
hypothesis of HD-FT21 on the complete exhibited factorization.
An empty tensor product is treated separately as one-dimensional.

\subsection*{5. A singleton standard layer: the precise complement replacement}

Suppose now that an actual faithful irreducible sphere-source
coefficient has the preceding complete quasiprimitive extraction,
and one \emph{whole singleton coarse layer} is ordinary
\(\operatorname{Std}_n\), with \(n\ge2^{30}\).
Its compatible natural quotient and the ordinary standard model
give genuine finite linear representations
\[
 V=(\operatorname{Std}_n\circ x)\otimes U,\qquad
 x:J\twoheadrightarrow T\in\{A_n,S_n\},\qquad
 \sum_i\operatorname{rank}(1-V(g_i))=2\dim V+6.
 \tag{QPX.10}
\]
All other original coarse layers form \(U\). Scalar phases
have been absorbed in \(U\). The quotient \(x\) is actual:
conjugation on the alternating core identifies its projective
normalizer with \(S_n\), and the ordinary standard matrices
give the faithful rational lift. No later core is independently
lifted to \(J\).

\begin{lemma}[Original-layer nonfaithfulness]
In (QPX.10), \(U\) is not faithful. Its kernel maps
injectively to \(T\) and has natural image containing \(A_n\).
\end{lemma}

\begin{proof}
Assume that \(U\) is faithful. A scalar \(U(g)\) makes \(g\)
central in \(J\), hence \(x(g)=1\). Thus \(x\) factors as
an actual surjection \(P_U\twoheadrightarrow T\).
If \(\dim U=1\), faithfulness already makes \(J\) abelian,
contrary to this quotient. Otherwise use the subdirect
embedding (QPX.8) in the \emph{remaining original} coarse
images. The exact coordinate-kernel lemma D-TD17 makes this
surjection factor through one coordinate \(P_j\twoheadrightarrow T\).
That group lemma only uses the actual subdirect embedding
and the normal subgroups of \(A_n,S_n\); it assumes no
primitivity of \(U\).

If \(P_j\) has Heisenberg core \(Q\), its abelian normal
image in \(T\) is trivial. The map therefore factors
through \(H=P_j/Q\). The Weyl centralizer calculation
identifies this quotient with its faithful symplectic action
on \(Q\), which is irreducible over the original
\(\mathbf F_p\) by Section~3. The original source, this
actual Heisenberg factor, the entire other partner, and
the surjections \(J\to P_j\to H\to T\) satisfy the standalone
D-TD43 hypotheses HAQ.1--2. That theorem gives a contradiction.
In particular its conclusion is not being applied merely
to an alternating composition factor of \(H\).

Otherwise the minimal normal core of \(P_j\) is \(N=S^t\).
It has trivial projective centralizer, so
\(N\le P_j\le\operatorname{Aut}(N)\).
Write \(\alpha:P_j\twoheadrightarrow T\).
If \(\alpha(N)\ne1\), it is the perfect normal subgroup
\(A_n\). The image of each simple direct factor is normal
in \(\alpha(N)\), hence trivial or \(A_n\). Transitivity
of \(P_j\) on the factors makes all of these images
nontrivial. Two copies of \(A_n\) cannot commute in
\(A_n\), so \(t=1\) and \(S\simeq A_n\).
The kernel of \(\alpha\) then has trivial intersection
with \(N\), so centralizes \(N\), and is trivial. Therefore
\(P_j\simeq T\), compatibly with the \emph{same} map \(x\).

If instead \(\alpha(N)=1\), the map factors through
\(P_j/N\le\operatorname{Out}(S)\wr S_t\).
The componentwise outer-automorphism kernel is solvable
by the same accepted Schreier input in D-TD17--18; its
normal image in \(T\) is trivial. Hence there is an actual
surjection from the transitive component-permutation image
to \(T\). The standalone D-TD45 excludes precisely this
actual tensor-orbit quotient at \(n\ge10^6\), retaining
the entire original partner. This rules out the second case.

We have reached a singleton quasisimple coordinate \(W\)
with projective image \(T\) and the same \(x\).
Genuine spin modules are excluded by the standalone
HD-FT28 at \(n\ge10^6\), and ordinary split modules by
HD-FT03 in its large-\(n\) actual-source scope.
For an ordinary extendible module choose an extension
\(S^\lambda\) to \(S_n\); its restriction to \(A_n\)
is irreducible and nontrivial. The ordinary branching
lemma proved in D-TD48 says that
\(\operatorname{Std}_n\otimes S^\lambda\) is reducible.
Every local \(W\)-operator is \(c_g S^\lambda(x(g))\):
both matrices implement the same automorphism of the
represented perfect \(A_n\), and their ratio is scalar
by Schur's lemma. A proper nonzero submodule of this
ordinary tensor product, tensored with the entire
remaining factor, is therefore \(J\)-invariant in \(V\).
This contradicts irreducibility.

Thus \(K_0=\ker U\ne1\). If \(x(k)=1\) for \(k\in K_0\),
then \(V(k)=I\); faithfulness gives \(k=1\).
The nontrivial normal subgroup \(x(K_0)\triangleleft T\)
contains \(A_n\), and is transitive in the natural action.
\end{proof}

This lemma replaces only the particular use of D-TD48
needed below. It does not apply D-TD12 afresh to an
arbitrary complement or assert that such a complement
is quasiprimitive.

\subsection*{6. The qualified minimal-source cutoff}

Retain the exact constant of D-TD54:
\[
 M_*=1+241\cdot240!\,2^{239},\qquad F_*=240M_*+2,\qquad
 N_*=[100(1+M_*)F_*^{240}]^2.
 \tag{QPX.11}
\]

\begin{theorem}[D-TD62]
Let an actual connected finite curve-cover family have a
faithful absolutely irreducible \emph{quasiprimitive}
coefficient representation with a model over its exact
imaginary quadratic character field \(K\). Suppose its
multiplicity variation has dimension six, signature \((3,3)\),
and actual period rank nine. Suppose that its naturally
polarized primitive sixfold, cut out by a rational
Rosati-fixed primitive coefficient projector, has nonsplit
signed determinant type.

Assume the coefficient degree is minimal among all actual
exact-\(K\) sources with this signature, actual period rank,
and the same natural signed determinant type. The comparison
class includes nonprimitive and nonquasiprimitive sources.
Fix any complete recursion of Sections~2--4.
Every alternating simple quotient in its quasisimple
extraction layers has natural degree strictly less than \(N_*\).
Its complete whole-layer complements consequently have
the uniform nonscalar support floor
\[
        \frac{\nu(S)}{\dim S}\ge\frac1{M_0},
        \qquad M_0=\max\{80,N_*-2\}.
 \tag{QPX.12}
\]
\end{theorem}

\begin{proof}
We give the full transfer through the accepted D-TD54
proof, identifying each place where its original
primitivity premise supplied a separate input.

Suppose that a largest alternating natural degree
\(n\ge N_*\) occurs in the chosen complete recursion.
Its actual source has degree greater than six: for an
ordinary alternating component use the elementary
Klein-four rank bound in HD-ML01, and for a genuine
spin component use the HD-FT23 degree floor.
The actual Euler and pointed-Hurwitz argument in D-BI18,
also written in full in D-TD56, gives a sphere base,
at least twelve genuine coefficient inertias, and the
budget in (QPX.10). This argument uses the actual
rank-nine variation and irreducibility, not primitivity.

The maximal-layer proof HD-ML01--02 transfers as follows.
The whole selected layer and its component/letter actions
are actual projective quotients by Sections~2--4.
Its entire original complement satisfies the floor
\(1/\max\{80,m-1\}\ge1/n\) by (QPX.9), since all its
alternating degrees \(m\) are at most \(n\).
There is no discarded or newly decomposed partner.
For a singleton ordinary layer, HD-FT03 excludes split
constituents, HD-FT16--17 exclude the higher ordinary
tails, and HD-FT22 excludes tail two by
\(4n^2\le61\max\{80,m-1\}\), which is impossible at
this maximal \(n\). For an ordinary layer with at
least two components, HD-TW05 excludes every extendible
tail at least two, HD-TW07 excludes split constituents,
and HD-AST14 excludes the standard tensor layer.
These theorems have the actual layer, full-partner floor,
and original-mark budget as hypotheses, not complex
primitivity of the original representation.
The unchanged integer (QPX.11) dominates every local
threshold used in HD-ML01.

For a spin layer the all-count exclusion HD-FT26 uses
the same original layer, all original powers, and the
same whole-partner floor. Its proof invokes the
standalone HD-FT25/23 estimates and contains no further
primitivity use. Thus it also transfers. We have an
actual \emph{singleton whole coarse layer}
\(\operatorname{Std}_n\), and (QPX.10) holds.

Put \(b=\dim U\). If \(b=1\), the actual scalar-standard
natural-polarization theorem D-BI08, with the exact
imaginary quadratic \(K\), forces the signed determinant
to be split. This contradicts the hypothesis, so \(b\ge2\).
Section~5 gives a nontrivial transitive \(K_0=\ker U\).
The exact-\(K\) image and model arguments HD-SS01/06
apply to this same representation and its letter-stabilizer
restriction; they require this actual kernel condition,
not primitivity.

If the whole-partner multiplicity \(\mu_U\) is zero,
the post-vanishing construction HD-SS06 supplies a
smaller actual degree-\(b\) sphere source, with the same
exact \(K\), signature, and actual period rank nine.
Its quotient correspondence and multiplicity-one
coefficient line scale the natural Hermitian pairing
by a positive rational number \(c\). Since the
multiplicity rank is six, \(c^6=\operatorname{Nm}_{K/\mathbf Q}(c^3)\);
the natural signed determinant type is unchanged.
The stated minimality excludes this source, proving
\(\mu_U>0\). No large-\(n/b\) assumption is needed in
this post-vanishing argument.

Apply the first, actual-source part of D-TD44 with
the complete remaining-original-layer floor (QPX.9).
Its complete-configuration application uses exactly
that floor. The additional observation there that \(U\)
is primitive is not needed: its original remaining
factor system supplies the HD-FT21 hypothesis directly.
Thus \(n\le108M_U-12\), where
\(M_U=\max\{80,m_{*,U}-1\}\). Since \(n\ge N_*>8628\),
there is an actual whole alternating layer in \(U\)
of degree \(m\) such that
\[
 \left\lceil\frac{n+120}{108}\right\rceil\le m\le n,
 \qquad m>241^{482}>2^{64}.
 \tag{QPX.13}
\]
The exact cutoff arithmetic is the unchanged D-TD49/51
calculation. The complete original support floor is at
least \(1/(n-1)\). The full-order tensor trace argument
of D-TD16, using that floor and the original twelve
genuine marks, bounds the transitive component count by
\[
       1\le t\le\left\lfloor\frac{n-1}{2}\right\rfloor<54m.
 \tag{QPX.14}
\]
The only use of original primitivity in D-TD16 was to
supply the complete floor. Its proof now has the same
input from (QPX.9).

The standalone estimates of D-TD49 exclude spin,
split and high-degree ordinary components in this
near-maximum layer; D-TD51 excludes all remaining
ordinary tails two through 240. Their near-maximum
applications use (QPX.13)--(QPX.14), the complete
original floor, actual normalized component actions,
and the original genuine-mark budget. These have all
just been supplied. Their character estimates and
induced-subset quotient constructions use no additional
original complex primitivity. Therefore the layer is
an ordinary standard layer of \(t\) components.

If \(t=1\), the two standard layer algebras are
separately stable. HD-MS03 uses their two actual
natural actions, the original marks, and the whole
remaining-partner floor \(1/n\) supplied by (QPX.9).
Its proof gives \(2m-1\le n\), while (QPX.13) gives
\(n<108m\). The proof of D-TD50 applies with the
following explicit replacement for its first step.
The original singleton \(n\)-layer has transitive
kernel \(K_0\) on its entire complement by Section~5.
For \(k\in K_0\), its ordinary \(m\)-standard matrix
is scalar, so its natural \(m\)-permutation is trivial
(\(m\ge5\)). The group is transitive on the \(m\)
letters, and \(K_0\) is transitive on the \(n\)
letters while fixing those \(m\) letters. Hence the
actual product action is transitive. This is precisely
the connected product-cover input in D-TD50.
Its subsequent full-power character comparison, all
finite and infinite mode cases, and the genus-zero
quadratic factorization use only this input and the
original budget. Primitivity in the final quadratic
factorization is that of the natural \(A_n/S_n\)
action, which is intrinsic and unchanged here.
Thus the two-standard configuration is impossible
at \(m>2^{64}\).

If \(t=2\), the standalone D-TD53 has exactly the
stable \(n\)-standard factor, actual transitive
two-component \(m\)-standard layer and entire remaining
partner. It has no original primitivity hypothesis,
and applies since \(n\ge m>2^{30}\).
If \(t\ge3\), the standalone D-TD52 applies to the
actual standard layer and arbitrary entire partner,
since \(m\ge2^{64}\) and \(3\le t\le54m\).
These exhaust all \(t\) and give the required contradiction.

The chosen recursion was arbitrary. Thus every one
has all alternating degrees below \(N_*\). Integer
degrees then give \(m-1\le N_*-2\), and (QPX.9)
proves (QPX.12).
\end{proof}

\subsection*{7. The enlargement does not identify quasiprimitivity with primitivity}

Here is an explicit exact-imaginary-quadratic, faithful,
nonmonomial example distinguishing the hypotheses at the
representation level. It is not claimed to supply a
rank-nine naturally nonsplit cover.

Let \(A_5\) act by conjugation on its six Sylow-five
subgroups. The deleted permutation character is
\[
       \psi_5=(5,1,-1,0,0)
       \quad\text{on }1,2A,3A,5A,5B.
 \tag{QPX.15}
\]
Indeed an involution fixes two Sylow-five subgroups:
their six order-ten normalizers contain thirty
involution incidences, and \(A_5\) has fifteen
involutions. A three-cycle fixes none; a five-cycle
fixes its unique Sylow-five subgroup. The character
norm is \((25+15+20)/60=1\), so this is a rational
irreducible module of dimension five. It is faithful
because \(A_5\) is simple. Let \(\psi_4\) be the
rational deleted natural five-letter module.
Its character \((4,0,1,-1,-1)\) has norm one as well.

Let \(\theta\) be either nontrivial character of
\(A_4/V_4\simeq C_3\). Restricting (QPX.15) to
\(A_4\) gives
\(\langle\operatorname{Res}\psi_5,\theta\rangle
=(5+3-4\zeta_3-4\zeta_3^2)/12=1\).
Dimension and Frobenius reciprocity therefore prove
\(\operatorname{Ind}_{A_4}^{A_5}\theta=\psi_5\).
For \(G=A_5\times A_5\times C_4\), with the faithful
character \(i\) of \(C_4\), put
\[
 \rho=\psi_5\otimes\psi_4\otimes i
 =\operatorname{Ind}_{A_4\times A_5\times C_4}^{G}
                    (\theta\otimes\psi_4\otimes i).
 \tag{QPX.16}
\]
This is a faithful absolutely irreducible degree-twenty
representation with an actual model and exact character
field \(K=\mathbf Q(i)\). Faithfulness follows because
a scalar matrix in either faithful \(A_5\)-factor
would represent a central element of that \(A_5\),
and hence the identity.

Every normal subgroup of \(G\) is a product of some
of its two \(A_5\)-factors with a subgroup of \(C_4\).
To see this, a nontrivial projection to either \(A_5\)
is onto, and commutation with that individual normal
factor puts the whole factor inside the subgroup.
After removing those factors the subgroup is central.
Restriction of \(\rho\) to any such product is
homogeneous. Thus \(\rho\) is quasiprimitive, whereas
(QPX.16) exhibits a nontrivial system of five
four-dimensional blocks.

It is not a monomial representation. If it were
induced from a linear character of an index-twenty
subgroup \(L_0\), the scalar group \(C_4\) would fix
every inducing line, and hence be contained in \(L_0\).
The image \(\overline L_0\subset A_5\times A_5\)
would have order \(180\).
Every proper subgroup of \(A_5\) has order at most
twelve: the only larger proper divisors \(15,20,30\)
would give faithful coset actions of degrees \(4,3,2\),
contradicting simplicity and \(|A_5|=60\).
Both projections of \(\overline L_0\) cannot be proper,
since then its order is at most \(144\).
If both are onto, the elementary subdirect-product
kernel argument for a simple group gives order
\(60\) or \(3600\), not \(180\).
If just one is onto, its kernel relative to the
other projection is normal in that full \(A_5\);
it cannot be trivial by the order bound, so it
contains that entire \(A_5\)-factor.
A linear inducing character is trivial on this
perfect factor, which is normal in \(G\), making
the induced representation trivial there. This
contradicts faithfulness of \(\rho\).

The displayed block system also illustrates why
D-TD60 does not automatically apply to every block
system. Its inducing subgroup is self-normalizing,
whereas the inducing character field is
\(\mathbf Q(i,\zeta_3)\), a proper quadratic
extension of \(K\). The nontrivial coefficient
Galois conjugate cannot be a normalizer conjugate,
so the block system is not Galois stable.
This is consistent with quasiprimitivity.

\paragraph{Exact boundary.}
D-TD62 supplies the complete extraction and the
qualified alternating cutoff for quasiprimitive
sources; it does not make those representations
primitive. The inducing-field obstruction in a
general nonhomogeneous source remains separate,
including proper-field cases left by D-TD56/60.
There is no new original group-order bound, general
induction-index bound, actual curve construction,
polarized-isogeny classification, or algebraic
cycle conclusion. The example in Section~7 is
representation-theoretic only. All actual geometric
consequences above retain the exact original field,
rank-nine variation, natural signed determinant,
and stated all-source minimality comparison.
